\section*{Capítulo \ref{ch:b2:proba} --- Probabilidad sobre espacios numerables}

\begin{solution}{exo:b2:proba:1}
Por simetría: el orden de extracción induce un orden relativo
uniformemente aleatorio sobre las bolas $1$ y $2$, de modo que $\P(1
\text{ antes que } 2) = \frac12$. Formalmente: intercambiar las
posiciones de las bolas $1$ y $2$ en una sucesión de extracciones es
una biyección de los resultados (equiprobables) que intercambia el
\omterm{def:b2:proba:space}{suceso} con su complementario. Entre las bolas $1, \dots, k$: el
orden relativo de esas $k$ bolas es uniforme entre las $k!$
ordenaciones, y la bola $1$ es la primera en $(k-1)!$ de ellas:
probabilidad $\frac{(k-1)!}{k!} = \frac1k$.
\end{solution}

\begin{solution}{exo:b2:proba:2}
$\frac{1}{k(k+1)} = \frac1k - \frac1{k+1}$, de modo que $\sum_{k\geq1}
p_k$ telescopa a $1$: es una \omterm{def:b2:proba:space}{medida de probabilidad}. Resultados
pares:
\[
\P(2\N^*) = \sum_{j=1}^{\infty}\frac{1}{2j(2j+1)}
= \sum_{j=1}^{\infty}\Bigl(\frac{1}{2j} - \frac{1}{2j+1}\Bigr)
= \frac12 - \frac13 + \frac14 - \frac15 + \cdots
\]
Esta es la serie armónica alternada sin su primer término y con los
signos cambiados: como $\ln 2 = 1 - \frac12 + \frac13 - \frac14 +
\cdots$ (\cref{ch:b2:series}),
\[
\P(2\N^*) = -\bigl(\ln 2 - 1\bigr) = 1 - \ln 2 \approx 0.307 .
\]
\end{solution}

\begin{solution}{exo:b2:proba:3}
Sean $S$ el \omterm{def:b2:proba:space}{suceso} de estar enfermo y $+$ el de dar positivo. Bayes
(\cref{thm:b2:proba:bayes}) con la partición $\{S, S^c\}$:
\[
\P(S \mid +)
= \frac{0.99 \times 10^{-4}}
       {0.99 \times 10^{-4} + 0.01 \times 0.9999}
= \frac{0.000099}{0.000099 + 0.009999}
\approx 0.0098 ,
\]
por debajo del $1\,\%$. Aunque el test tenga un ``$99\,\%$ de
acierto'', un resultado positivo deja al paciente con
aproximadamente un $99\,\%$ de probabilidad de estar sano: los falsos
positivos de la inmensa mayoría sana desbordan a los verdaderos
positivos de la minúscula minoría enferma. Los tests de cribado para
enfermedades raras han de leerse siempre a través de este cálculo de
tasa base.
\end{solution}

\begin{solution}{exo:b2:proba:4}
Punto a punto sobre $\Omega$: $\omega \in \bigcup A_i$ si y solo si
se anula algún factor $1 - \mathbf{1}_{A_i}(\omega)$, de modo que
\[
\mathbf{1}_{\bigcup A_i}
= 1 - \prod_{i=1}^n\bigl(1 - \mathbf{1}_{A_i}\bigr)
= \sum_{\emptyset \neq J \subseteq \{1,\dots,n\}}
(-1)^{\abs J + 1}\prod_{i \in J}\mathbf{1}_{A_i} ,
\]
desarrollando el producto y pasando el $1$ al otro lado. Ahora bien,
$\prod_{i\in J}\mathbf{1}_{A_i} = \mathbf{1}_{\bigcap_{i \in J}
A_i}$, y sumar contra los pesos $\P(\{\omega\})$ ---legítimo:
finitos términos acotados, cada familia sumable--- convierte cada
indicatriz en la probabilidad de su \omterm{def:b2:proba:space}{suceso}, lo que da la fórmula.
\end{solution}

\begin{solution}{exo:b2:proba:5}
Sea $A_i$ el \omterm{def:b2:proba:space}{suceso} ``la carta $i$ está en el sobre correcto''. Para
$J$ de tamaño $k$, $\P\bigl(\bigcap_{i\in J}A_i\bigr) =
\frac{(n-k)!}{n!}$ (fíjense $k$ cartas y permútense las demás). Por
inclusión-exclusión,
\[
\P\Bigl(\bigcup A_i\Bigr)
= \sum_{k=1}^n (-1)^{k+1}\binom nk \frac{(n-k)!}{n!}
= \sum_{k=1}^n \frac{(-1)^{k+1}}{k!} ,
\]
de modo que
\[
\P(\text{ningún acierto})
= 1 - \P\Bigl(\bigcup A_i\Bigr)
= \sum_{k=0}^{n}\frac{(-1)^k}{k!}
\xrightarrow[n\to\infty]{} e^{-1} \approx 0.368 .
\]
Exactamente un acierto: una permutación con exactamente un punto
fijo queda determinada por la elección de la carta fija ($n$ maneras)
y un \emph{desarreglo} (una disposición sin aciertos) de las otras $n
- 1$; escribiendo $D_{n-1} = (n-1)!\sum_{k=0}^{n-1}\frac{(-1)^k}{k!}$
para el número de desarreglos (la primera parte, escalada por
$(n-1)!$),
\[
\P(\text{exactamente un acierto})
= \frac{n\,D_{n-1}}{n!}
= \frac{D_{n-1}}{(n-1)!}
= \sum_{k=0}^{n-1}\frac{(-1)^k}{k!}
\xrightarrow[n\to\infty]{} e^{-1} :
\]
en el límite, ``ningún acierto'' y ``exactamente un acierto'' son
igual de probables, cada uno con probabilidad $e^{-1}$.
\end{solution}

\begin{solution}{exo:b2:proba:6}
Condiciónese al comienzo (regla de la cadena,
\cref{thm:b2:proba:bayes}):
\begin{itemize}
\item primer lanzamiento cruz (probabilidad $1 - p$): el juego
recomienza de cero, y durar más de $n$ significa durar más de $n - 1$
desde ahí: contribución $(1-p)\,q_{n-1}$;
\item primeros lanzamientos cara-cruz (probabilidad $p(1-p)$): se
recomienza tras dos lanzamientos: contribución $p(1-p)\,q_{n-2}$;
\item primeros lanzamientos cara-cara: el juego ha terminado (dentro
de $n$ lanzamientos, con $n \geq 2$): contribución $0$.
\end{itemize}
De ahí, $q_n = (1-p)q_{n-1} + p(1-p)q_{n-2}$. La ecuación
característica $r^2 = (1-p)r + p(1-p)$ tiene raíces
\[
r_\pm = \frac{(1-p) \pm \sqrt{(1-p)^2 + 4p(1-p)}}{2},
\]
con $\abs{r_\pm} < 1$: en efecto, el polinomio $\chi(r) = r^2 -
(1-p)r - p(1-p)$ cumple $\chi(1) = 1 - (1-p) - p(1-p) = p^2 > 0$ y
$\chi(-1) = 1 + (1-p) - p(1-p) > 0$, mientras que $\chi(0) = -p(1-p)
< 0$: una raíz en $\intoo{-1}{0}$ y otra en $\intoo{0}{1}$. Así pues,
$q_n = \alpha r_+^n + \beta r_-^n \to 0$. Los \omterm{def:b2:proba:space}{sucesos} ``el juego dura
más de $n$'' decrecen hacia ``el juego no termina nunca''; la
\omterm{def:b2:metric:continuity}{continuidad} monótona (\cref{thm:b2:proba:continuity}) da $\P(\text{no
termina nunca}) = \lim q_n = 0$: el juego termina casi seguramente.
\end{solution}

\begin{solution}{exo:b2:proba:7}
\emph{$\P(R_n) = 1/n$:} entre las $n$ primeras extracciones, cada una
de las $n$ posiciones relativas de la última es igual de probable
(uniformidad del orden relativo), y $R_n$ es el \omterm{def:b2:proba:space}{suceso} de que sea la
mayor: probabilidad $1/n$.

\emph{Infinitos récords:} $\sum_n \P(R_n) = \sum 1/n = \infty$ y los
$R_n$ son \omterm{def:b2:proba:independence}{independientes} (admitido), así que Borel--Cantelli~2
(\cref{thm:b2:proba:borelcantelli}) da $\P(\limsup R_n) = 1$: los
récords no cesan nunca, casi seguramente; pero se van rarificando
logarítmicamente.

\emph{Récords consecutivos:} por \omterm{def:b2:proba:independence}{independencia},
\[
\sum_n \P(R_n \cap R_{n+1})
= \sum_n \frac{1}{n(n+1)} < \infty ,
\]
de modo que se aplica Borel--Cantelli~1: casi seguramente, solo un
número finito de veces un récord va seguido inmediatamente de otro
récord. Las dos mitades del lema trabajan en tándem: infinitos
récords, pero (c.s.) a partir de cierto punto nunca dos seguidos.
\end{solution}

\begin{solution}{exo:b2:proba:8}
\emph{Primera parte:} $\sum \P(A_n) = \sum\frac{1}{n+1} = \infty$ con
\omterm{def:b2:proba:independence}{independencia}: Borel--Cantelli~2 da $\P(\limsup A_n) = 1$. Cada $A_n$
individual es cada vez más improbable y, aun así, casi todo $\omega$
pertenece a infinitos de ellos.

\emph{Contraejemplo sin \omterm{def:b2:proba:independence}{independencia}:} tómese $\Omega = \N^*$ con los
pesos $p_k = \frac{1}{k(k+1)}$ del~\cref{exo:b2:proba:2}, y $A_n =
\{k \in \N^* : k \geq n\}$. Entonces
\[
\P(A_n) = \sum_{k \geq n}\Bigl(\frac1k - \frac1{k+1}\Bigr)
= \frac1n ,
\qquad
\sum_n \P(A_n) = \infty ,
\]
pero los $A_n$ son decrecientes, de modo que $\limsup_n A_n =
\bigcap_n A_n = \emptyset$: $\P(\limsup A_n) = 0$. La divergencia de
$\sum\P(A_n)$ por sí sola no garantiza nada cuando los \omterm{def:b2:proba:space}{sucesos} se
amontonan sobre una parte cada vez más pequeña del espacio; la
\omterm{def:b2:proba:independence}{independencia} es lo que prohíbe esa conspiración.
\end{solution}

\begin{solution}{exo:b2:proba:9}
Con $q = 1 - p$, la primera cara cae en el rango $2j + 1$ con
probabilidad $q^{2j}p$, de modo que
\[
\P(\text{rango impar}) = \sum_{j\geq0}q^{2j}p
= \frac{p}{1 - q^2} = \frac{1}{1 + q} .
\]
Para una moneda equilibrada: $\frac1{1 + 1/2} = \frac23$.
(Comprobación de sensatez: los rangos impares deberían ser más
probables, ya que el rango $1$ va primero; y en efecto, $\frac1{1+q}
> \frac12$ siempre.)
\end{solution}

\begin{solution}{exo:b2:proba:10}
Los \omterm{def:b2:proba:space}{sucesos} $B_N = \bigcap_{n=1}^N A_n^c$ decrecen hacia
$\bigcap_nA_n^c$ y, por la \omterm{def:b2:proba:independence}{independencia} de los complementarios,
$\P(B_N) = \prod_{n=1}^N(1 - p_n)$; la \omterm{def:b2:metric:continuity}{continuidad} monótona
(\cref{thm:b2:proba:continuity}) da el límite mostrado. Tomando
logaritmos, $\prod(1 - p_n) > 0$ si y solo si $\sum-\ln(1 - p_n) <
\infty$. Si $\sum p_n < \infty$, entonces $p_n \to 0$ y $-\ln(1 -
p_n) \sim p_n$: la serie de logaritmos converge. Si $\sum p_n =
\infty$, entonces $-\ln(1 - p_n) \geq p_n$ fuerza la divergencia, y
el producto es $0$. Esto casa con Borel--Cantelli~2: para $\sum p_n =
\infty$, no solo $\P(\text{no ocurre ningún }A_n) = 0$, sino que casi
seguramente ocurren infinitos $A_n$.
\end{solution}

\begin{solution}{exo:b2:proba:11}
Digamos que la caja $A$ es la primera que se descubre vacía, y que la
otra guarda $k$ cerillas. Esto significa que, entre las primeras $2n
- k$ veces que mete la mano, exactamente $n$ fueron a $A$ y $n - k$ a
$B$ (en algún orden), y que la vez número $2n - k + 1$ fue de nuevo a
$A$, encontrándola vacía. Las metidas de mano son elecciones
equilibradas \omterm{def:b2:proba:independence}{independientes}, de modo que este \omterm{def:b2:proba:space}{suceso} tiene
probabilidad $\binom{2n-k}{n}2^{-(2n-k)}\cdot\frac12$; y duplicando
(la caja vacía puede ser cualquiera de las dos) se obtiene
\[
\P(\text{la otra caja tiene }k) = \binom{2n-k}{n}\,2^{-(2n-k)} .
\]
Para $n = 1$: $k = 1$ da $\binom11 2^{-1} = \frac12$ y $k = 0$ da
$\binom21 2^{-2} = \frac12$: total $1$, como debe ser.
\end{solution}

\begin{solution}{exo:b2:proba:12}
(a) Si $\P(\{n\}) = c$ para todo $n$, la $\sigma$-aditividad fuerza
$1 = \sum_nc$: imposible, tanto si $c = 0$ (suma $0$) como si $c > 0$
(suma infinita). No hay probabilidad uniforme sobre $\N$.

(b) Si $A \cap B = \emptyset$ y existen $d(A)$ y $d(B)$, entonces
$\abs{(A \sqcup B)\cap\intint1n} = \abs{A\cap\intint1n} +
\abs{B\cap\intint1n}$, de donde $d(A \sqcup B) = d(A) + d(B)$:
aditividad finita sobre tales parejas. Todo conjunto unitario tiene
función de conteo finalmente constante, luego densidad $0$, mientras
que $d(\N^*) = 1$. Si $d$ fuera $\sigma$-aditiva, $\N^* =
\bigsqcup_k\{k\}$ daría $1 = \sum_k 0 = 0$: la densidad es
finitamente aditiva, pero no $\sigma$-aditiva; el axioma tiene
contenido.

(c) Sea $A = \bigcup_{k\geq0}\intint{4^k}{2\cdot4^k - 1}$ (bloques de
$4^k$ a $2\cdot4^k - 1$). En $n = 2\cdot4^K - 1$, el recuento es
$\sum_{k\leq K}4^k \sim \frac43 4^K$, lo que da razón $\to \frac23$;
en $n = 4^{K+1} - 1$, el recuento no ha cambiado, lo que da razón
$\to \frac13$. La razón oscila entre los límites $\frac13$ y
$\frac23$: no hay densidad.
\end{solution}

\begin{solution}{pb:b2:proba:1}
\textbf{1.} Un camino de longitud $n$ queda determinado por el
conjunto de sus pasos hacia arriba; terminar en $k$ significa $u$
pasos arriba y $n - u$ abajo con $u - (n - u) = k$, es decir, $u =
\frac{n+k}2$: posible si y solo si $n + k$ es par y $\abs k \leq n$,
de $\binom{n}{(n+k)/2}$ maneras. Cada camino concreto es un punto de
la medida producto equilibrada sobre $n$ lanzamientos: probabilidad
$2^{-n}$. De ahí, $\P(S_n = k) = N_n(k)2^{-n}$.

\textbf{2.} $S_n$ tiene la paridad de $n$, luego $S_{2n+1} \neq 0$; y
$u_n = N_{2n}(0)4^{-n} = \binom{2n}n4^{-n}$. Valores: $u_1 =
\frac12$, $u_2 = \frac6{16} = \frac38$, $u_3 = \frac{20}{64} =
\frac5{16}$.

\textbf{3.} $\dfrac{u_n}{u_{n-1}} =
\dfrac{\binom{2n}n}{4\binom{2n-2}{n-1}} = \dfrac{(2n)(2n-1)}{4n^2} =
\dfrac{2n-1}{2n} < 1$: decreciente. Por el~\cref{ex:b2:comparison:centralbinomial}, $\binom{2n}n \sim
\frac{4^n}{\sqrt{\pi n}}$, luego $u_n \sim \frac1{\sqrt{\pi n}} \to
0$, y $\sum u_n$ diverge por comparación con $\sum n^{-1/2}$.

\textbf{4.} Dado un camino de $1$ a $k$ que toca $0$, refléjese su
segmento inicial (hasta la \emph{primera} visita a $0$) respecto del
eje horizontal: el resultado es un camino de $-1$ a $k$, y la
operación es una involución; todo camino de $-1$ a $k \geq 1$ ha de
cruzar $0$, y reflejar de vuelta su segmento inicial recupera el
original. Por tanto, los caminos que tocan son $N_{n-1}(k + 1)$
(pues de $-1$ a $k$ el desplazamiento es $k + 1$). Un camino de $0$ a
$k$ que permanece $> 0$ tras el instante $0$ empieza con un paso
hacia arriba y luego va de $1$ a $k$ en $n - 1$ pasos sin tocar $0$:
hay $N_{n-1}(k-1) - N_{n-1}(k+1)$ de ellos.

\textbf{5.} Con $m = \frac{n+k}2$, usando $\binom{n-1}{m-1} = \frac
mn\binom nm$ y $\binom{n-1}{m} = \frac{n-m}n\binom nm$:
\[
\frac{N_{n-1}(k-1) - N_{n-1}(k+1)}{N_n(k)}
= \frac{\binom{n-1}{m-1} - \binom{n-1}{m}}{\binom nm}
= \frac{m - (n - m)}{n} = \frac kn .
\]
Para $n = 3$, $k = 1$: $N_3(1) = 3$ caminos ($++-$, $+-+$, $-++$), de
los cuales solo $++-$ permanece positivo ($+-+$ vuelve a $0$ en el
instante $2$): uno de tres, y $\frac kn = \frac13$.

\textbf{6.} Por simetría, la probabilidad es $2\P(S_i > 0\ \forall i
\leq 2n)$. Sumando sobre el punto final $2k$ y usando la pregunta 4
(con $n$ sustituido por $2n$):
\[
\P(S_i > 0\ \forall i) = 2^{-2n}\sum_{k\geq1}
\bigl(N_{2n-1}(2k-1) - N_{2n-1}(2k+1)\bigr)
= 2^{-2n}\,N_{2n-1}(1),
\]
una suma telescópica. Ahora bien, $N_{2n-1}(1) = \binom{2n-1}{n}$ y
$2\binom{2n-1}n = \binom{2n}n$ (Pascal), de modo que la probabilidad
mostrada es $2\cdot2^{-2n}\binom{2n-1}n = \binom{2n}n4^{-n} = u_n$.

\textbf{7.} Los \omterm{def:b2:proba:space}{sucesos} $D_n = \{S_i \neq 0,\ i \leq 2n\}$
decrecen, con intersección ``no hay retorno nunca''; por la
\omterm{def:b2:metric:continuity}{continuidad} monótona y la pregunta 6, $\P(\text{sin retorno}) = \lim
u_n = 0$: el paseo vuelve casi seguramente. Además, $f_n =
\P(D_{n-1}) - \P(D_n) = u_{n-1} - u_n$ y, por la pregunta 3,
\[
u_{n-1} - u_n = u_n\Bigl(\frac{2n}{2n-1} - 1\Bigr) =
\frac{u_n}{2n-1};
\qquad
\sum_{n\geq1}f_n = u_0 - \lim u_n = 1 .
\]

\textbf{8.} $2n\,f_n = \frac{2n}{2n-1}u_n \geq u_n$, y $\sum u_n =
\infty$ (pregunta 3): la serie $\sum 2nf_n$ diverge. El primer
retorno es cierto, pero no tiene tiempo medio de espera finito; el
paseo es \emph{recurrente nulo}, en el vocabulario que proporcionará
el~\cref{ch:b2:randomvar}.

\textbf{9.} El \omterm{def:b2:proba:space}{suceso} ``al menos $k$ retornos'' es la unión
\omterm{def:b2:structures:countable}{numerable} disjunta, sobre $0 < n_1 < \dots < n_k$, de los \omterm{def:b2:proba:space}{sucesos}
``los $k$ primeros retornos ocurren exactamente en los instantes
$2n_1, \dots, 2n_k$''. Un \omterm{def:b2:proba:space}{suceso} así es la intersección de $k$
\omterm{def:b2:proba:space}{sucesos} que dependen de los bloques disjuntos de lanzamientos
$\intint1{2n_1}$, $\intint{2n_1+1}{2n_2}$, \dots, exigiendo cada
bloque que un paseo nuevo haga su primer retorno tras exactamente el
número asignado de pasos; por la \omterm{def:b2:proba:independence}{independencia} de los bloques, su
probabilidad es $f_{n_1}f_{n_2-n_1}\cdots f_{n_k-n_{k-1}}$. Sumando
por paquetes (\cref{ch:b2:series}, con todos los términos no
negativos):
\[
\P(\text{al menos }k\text{ retornos})
= \Bigl(\sum_{n\geq1}f_n\Bigr)^{\!k} = 1^k = 1 .
\]
Los \omterm{def:b2:proba:space}{sucesos} decrecen en $k$, de modo que, por la \omterm{def:b2:metric:continuity}{continuidad}
monótona, $\P(\text{infinitos retornos}) = 1$: recurrencia.

\textbf{10.} Por la pregunta 9, el paseo hace infinitas excursiones
lejos de $0$. El primer paso de cada excursión es una moneda nueva,
\omterm{def:b2:proba:independence}{independiente} de todo lo anterior: la probabilidad de que las $m$
primeras excursiones empiecen todas hacia abajo es $2^{-m}$. Para
alcanzar $1$, al paseo le basta con un inicio de excursión hacia
arriba (desde $<0$ ha de pasar por $0$ antes de llegar a $1$, pues
los pasos son $\pm1$), luego $\P(\text{no alcanzar nunca }1) \leq
2^{-m}$ para todo $m$: el paseo alcanza $1$ casi seguramente.
Descomponiendo sobre el instante de llegada (casi seguramente
finito), el paseo reiniciado ahí es un paseo nuevo que arranca en
$1$: por inducción alcanza todo $k \geq 1$ casi seguramente y, por
simetría, todo $k \leq -1$. Por último, reiniciando en la primera
visita a $k$, la pregunta 9 se aplica al paseo nuevo: todo sitio se
visita infinitas veces, casi seguramente.

\textbf{11.} Los \omterm{def:b2:proba:space}{sucesos} $A_n = \{S_{2n} = 0\}$ distan mucho de ser
\omterm{def:b2:proba:independence}{independientes} (estar en $0$ en el instante $2n$ hace mucho más
probable estar en $0$ en el instante $2n + 2$ de lo que indica
$u_{n+1}$), de modo que Borel--Cantelli~2 no está disponible; y, en
efecto, todo el trabajo de la parte II consistió en sustituirlo. La
otra dirección no necesita \omterm{def:b2:proba:independence}{independencia}: \emph{si} $\sum\P(A_n)$
converge, Borel--Cantelli~1 da un número finito de retornos casi
seguramente. Esa implicación es el motor de toda demostración de
transitoriedad de más abajo.

\textbf{12.} Un retorno en el instante $2n$ exige $n$ pasos arriba y
$n$ abajo: $\P(S_{2n} = 0) = \binom{2n}np^nq^n = u_n(4pq)^n$, y $4pq
= 1 - (p - q)^2 < 1$ para $p \neq \frac12$. Como $u_n \leq 1$, la
serie $\sum\P(S_{2n} = 0)$ está dominada por la geométrica
$\sum(4pq)^n$: convergente. Por Borel--Cantelli~1, $\P(S_{2n} = 0
\text{ infinitas veces}) = 0$: un número finito de retornos, casi
seguramente.

\textbf{13.} Para $n + k$ par, $\P(S_n = k) =
\binom{n}{\frac{n+k}2}p^{\frac{n+k}2}q^{\frac{n-k}2}$; el coeficiente
binomial es a lo sumo el central, y $p^{\frac{n+k}2}q^{\frac{n-k}2} =
(pq)^{n/2}(p/q)^{k/2}$, lo que da la cota enunciada $\leq
2^n(pq)^{n/2}(p/q)^{k/2} = (4pq)^{n/2}(p/q)^{k/2}$, \omterm{def:b2:series:summable}{sumable} en $n$
porque $\sqrt{4pq} < 1$. Borel--Cantelli~1: el sitio $k$ se visita un
número finito de veces casi seguramente; y la unión sobre $k \in \Z$
de los \omterm{def:b2:proba:space}{sucesos} nulos excepcionales sigue siendo nula (subaditividad
\omterm{def:b2:structures:countable}{numerable}). Casi seguramente, todo sitio se visita un número finito
de veces, de modo que la sucesión de enteros $(S_n)$ abandona
definitivamente toda ventana acotada: $\abs{S_n} \to \infty$.

\textbf{14.} $\P(S_n \neq 0,\ 1 \leq n \leq 200) = u_{100} =
\binom{200}{100}4^{-100} \approx \frac1{\sqrt{100\pi}} \approx
0.056$: más de una probabilidad entre veinte de que $200$
lanzamientos equilibrados no empaten nunca. El decaimiento
$1/\sqrt{\pi n}$ es dolorosamente lento: la certeza de un empate
(pregunta 7) es compatible con tramos larguísimos sin empates; un
primer sabor de los fenómenos del arcoseno de la parte IV.

\textbf{15.} Condiciónese al primer paso. Si $X_1 = +1$, entonces
$T_1 = 1$, y $f_1 = \frac12$ concuerda. Si $X_1 = -1$, el paseo ha de
subir de $-1$ a $1$; por la descomposición en bloques, volver a $0$
por primera vez en el instante $2n$ se escinde en: un paso abajo y
después un paseo nuevo que arranca en $-1$ y alcanza $0$ por primera
vez ---equivalentemente, un paseo nuevo que alcanza $+1$ por primera
vez--- en $2n - 1$ pasos, o el suceso simétrico hacia arriba. Ambos
signos contribuyen por igual:
\[
f_n = 2\cdot\tfrac12\,\P(T_1 = 2n - 1) = \P(T_1 = 2n-1) .
\]
De ahí, $\P(T_1 < \infty) = \sum f_n = 1$, mientras que $\sum_n(2n -
1)f_n = \sum_n u_n = \infty$ por la pregunta 7: el paseo alcanza $1$
casi seguramente, en tiempo medio infinito.

\textbf{16.} Pártase $\{M_n \geq k\}$ según el valor final $S_n = m$.
Para $m \geq k$, la condición $M_n \geq k$ es automática. Para $m <
k$, refléjese el camino tras su \emph{primera} visita al nivel $k$:
esto es una biyección entre $\{M_n \geq k, S_n = m\}$ y $\{S_n = 2k -
m\}$ (todo camino que termina en $2k - m > k$ visita $k$; y reflejar
de vuelta es la inversa). De ahí,
\[
\P(M_n \geq k)
= \sum_{m > k}\P(S_n = m) + \P(S_n = k)
+ \sum_{m < k}\P(S_n = 2k - m)
= 2\P(S_n > k) + \P(S_n = k).
\]

\textbf{17.} En el instante par $2n$ con $k = 1$: $\P(S_{2n} = 1) =
0$ y $\P(S_{2n} > 1) = \P(S_{2n} \geq 2)$, de modo que
\[
\P(M_{2n} \geq 1) = 2\P(S_{2n} \geq 2)
= \P(S_{2n} \geq 2) + \P(S_{2n} \leq -2)
= 1 - u_n .
\]
Por tanto, $\P(S_i \leq 0\ \forall i \leq 2n) = u_n$: el paseo no va
por delante en los $2n$ primeros pasos exactamente con la misma
frecuencia con que no empata (pregunta 6); dos \omterm{def:b2:proba:space}{sucesos} bastante
distintos, llevados por el mismo $u_n$.

\textbf{18.} $\{L_{2n} = 2k\} = \{S_{2k} = 0\} \cap \{\text{el paseo
de los lanzamientos } 2k+1, \dots, 2n \text{ no tiene ceros}\}$. Los
dos \omterm{def:b2:proba:space}{sucesos} dependen de bloques disjuntos de lanzamientos, así que
son \omterm{def:b2:proba:independence}{independientes}; el primero tiene probabilidad $u_k$, y el
segundo $u_{n-k}$ por la pregunta 6 aplicada al paseo nuevo de $2n -
2k$ pasos. De ahí, $\P(L_{2n} = 2k) = u_ku_{n-k}$. Como $L_{2n}$ toma
exactamente los valores $0, 2, \dots, 2n$, esas probabilidades suman
$1$: $\sum_{k=0}^nu_ku_{n-k} = 1$, una identidad binomial entregada
por una partición probabilista.

\textbf{19.} La simetría es inmediata: $u_ku_{n-k} = u_{n-k}u_k$.
Como $u_j$ decrece en $j$, el producto $u_ku_{n-k}$ es mínimo para
$k$ central y máximo en los extremos $k \in \{0, n\}$, donde vale
$u_n$; cuantitativamente, $u_ku_{n-k} \approx
\frac1{\pi\sqrt{k(n-k)}}$ en el grueso, frente a $u_n \approx
\frac1{\sqrt{\pi n}}$ en los bordes. Para $n = 5$: $\P(L_{10} = 0) =
\P(L_{10} = 10) = u_5 = \frac{63}{256} \approx 0.246$, mientras que
$\P(L_{10} = 4) = u_2u_3 = \frac38\cdot\frac5{16} = \frac{15}{128}
\approx 0.117$. En una partida equilibrada larga, la última
igualación es más probable muy al principio o muy al final: un
jugador suele ir por delante durante tramos enormes, sin ningún sesgo
en la moneda.

\textbf{20.} La imagen: en el instante $n$ el paseo vive a escala
$\sqrt n$ (la dispersión binomial de la pregunta 3; $u_n \sim
1/\sqrt{\pi n}$ es la altura del pico central); vuelve a $0$
infinitas veces con probabilidad $1$ (parte II) y, sin embargo, el
tiempo de espera entre retornos tiene media divergente (pregunta 8),
razón por la cual una sola excursión puede ocupar una fracción
positiva de cualquier horizonte; en consonancia, el último empate de
una partida de $2n$ pasos queda repartido con los valores extremos
como los más probables (preguntas 18--19), y no ir nunca por delante
tiene la misma probabilidad de decaimiento lento $u_n$ que no empatar
nunca (pregunta 17). Certeza en el límite, persistencia en todo
horizonte finito: ese es el paseo equilibrado.

\textbf{21.} Pártase $\{S_{2n} = 0\}$ ($n \geq 1$) según el instante
del primer retorno $2k$, $1 \leq k \leq n$: el primer bloque de $2k$
lanzamientos realiza un primer retorno, los $2n - 2k$ lanzamientos
restantes realizan un retorno de un paseo nuevo, y los bloques son
\omterm{def:b2:proba:independence}{independientes}: $u_n = \sum_{k=1}^nf_ku_{n-k}$. Ambas series $U(x)
= \sum u_nx^n$ y $F(x) = \sum f_nx^n$ tienen radio $\geq 1$
(coeficientes en $\intcc01$), y el \omterm{thm:b2:series:fubini}{producto de Cauchy}
(\cref{ch:b2:powerseries}) da, para $0 \leq x < 1$,
\[
U(x) - 1 = \sum_{n\geq1}\Bigl(\sum_{k=1}^n
f_ku_{n-k}\Bigr)x^n = F(x)\,U(x),
\qquad\text{es decir,}\qquad
U(x)\bigl(1 - F(x)\bigr) = 1 .
\]

\textbf{22.} Cuando $x \uparrow 1$, $U(x)$ y $F(x)$ crecen
(coeficientes no negativos); toda suma parcial $\sum_{n\leq N}u_n$ es
límite de $\sum_{n\leq N}u_nx^n \leq U(x)$, de modo que $U(x)
\uparrow \sum u_n \in \intoc0{+\infty}$, e igualmente $F(x) \uparrow
f = \sum f_n$. Si $\sum u_n = \infty$: $1 - F(x) = 1/U(x) \to 0$,
luego $f = 1$. Si $\sum u_n = S < \infty$: $1 - f = 1/S > 0$, luego
$f < 1$. Comprobaciones: paseo equilibrado, $\sum u_n = \infty$ y $f
= 1$ (preguntas 3 y 7); paseo sesgado, $\sum u_n(4pq)^n < \infty$ y,
en consonancia, $f = 1 - 1/\sum_{n\geq0}u_n(4pq)^n < 1$, coherente
con la finitud casi segura del número de retornos (pregunta 12).

\textbf{23.} Para los cuatro pasos $(\pm1, 0), (0, \pm1)$ del paseo
sobre $\Z^2$, los incrementos de $U = X + Y$ y $V = X - Y$ son:
$(+,+)$ para $(1,0)$, $(+,-)$ para $(0,1)$, $(-,+)$ para $(0,-1)$ y
$(-,-)$ para $(-1,0)$; cada pareja de signos con probabilidad
$\frac14 = \frac12\cdot\frac12$: los dos paseos coordenados $(U_n)$ y
$(V_n)$ son paseos equilibrados \omterm{def:b2:proba:independence}{independientes} sobre $\Z$. Como
$S^{(2)}_{2n} = (0,0)$ si y solo si $U_{2n} = 0$ y $V_{2n} = 0$,
\[
\P\bigl(S^{(2)}_{2n} = (0,0)\bigr) = u_n^2 \sim \frac1{\pi
n}, \qquad \sum_nu_n^2 = \infty .
\]
La identidad de renovación de la pregunta 21 y la dicotomía de la
pregunta 22 no usaron nada unidimensional (solo la descomposición
sobre el primer retorno y la \omterm{def:b2:proba:independence}{independencia} de bloques disjuntos), de
modo que $\sum u_n^{(2)} = \infty$ da $f^{(2)} = 1$, y el argumento
de la pregunta 9 lo mejora: el paseo sobre $\Z^2$ vuelve al origen
infinitas veces casi seguramente.

\textbf{24.} Con la cota admitida $\P(S^{(3)}_{2n} = 0) \leq
Cn^{-3/2}$, la serie converge, y Borel--Cantelli~1 da un número
finito de retornos casi seguramente: el paseo sobre $\Z^3$ es
transitorio (y la misma cota con exponente $-d/2$ cubre todo $d \geq
3$). En conjunto: el \emph{teorema de Pólya}: el
\omterm{pb:b2:proba:1}{paseo aleatorio simple} es recurrente sobre $\Z$ y $\Z^2$, y
transitorio sobre $\Z^d$ para $d \geq 3$. Un borracho encuentra el
camino a casa; un pájaro borracho puede que no.

\textbf{25.} El recuento de caminos y la reflexión produjeron las
leyes exactas ($u_n$, el teorema de la papeleta, $f_n$, el máximo, el
último cero); la \omterm{def:b2:metric:continuity}{continuidad} monótona convirtió todo enunciado
límite (``vuelve al menos una vez'', ``infinitas veces'') en un
límite de probabilidades de horizonte finito; la \omterm{def:b2:proba:independence}{independencia} de
bloques disjuntos impulsó las descomposiciones de renovación
(preguntas 9, 18, 21), y es el esqueleto \omterm{def:b2:structures:countable}{numerable} de la propiedad de
Markov; Borel--Cantelli~1 fue el arma de la transitoriedad (preguntas
12--13 y 24), sin necesitar \omterm{def:b2:proba:independence}{independencia}; y la identidad de
renovación lo organizó todo en la dicotomía $\sum u_n = \infty \iff$
recurrencia. El único insumo \omterm{def:b2:powerseries:analytic}{analítico} es la estimación local $u_n
\sim 1/\sqrt{\pi n}$: su cuadrado $1/(\pi n)$ sigue divergiendo
(dimensión $2$, recurrente), mientras que $n^{-3/2}$ converge
(dimensión $3$, transitorio); el teorema de Pólya es, al final, un
enunciado sobre la divergencia de $\sum n^{-d/2}$.
\end{solution}
